\title{LongRoPE: Extending LLM Context Window Beyond 2 Million Tokens}

\begin{abstract}
Expanding the context window in large language models (LLMs) is highly sought after. Nevertheless, current models only achieve context windows of roughly 128k tokens, largely constrained by expensive fine-tuning procedures, limited availability of extended textual data, and problematic values caused by unfamiliar token positions.

This study presents LongRoPE, a novel approach that extends the context window of pre-trained LLMs to a remarkable \textbf{2048k} tokens—marking the first time such length has been achieved—using at most 1k fine-tuning steps restricted to 256k token sequences, all while preserving the performance levels observed at the native short context window. The success of LongRoPE hinges on three main methodological advances: (i) identification and utilization of two distinct types of non-uniformities in positional interpolation through an efficient search algorithm, yielding better initialization for fine-tuning and facilitating up to an 8$\times$ window extension in zero-shot scenarios; (ii) a staged extension technique, wherein a model is initially fine-tuned to handle 256k tokens and subsequently undergoes a second round of positional interpolation, enabling the final 2048k-token capability; and (iii) refinement of LongRoPE on 8k-token sequences to restore high accuracy for shorter contexts. Comprehensive evaluations conducted with LLaMA2 and Mistral across multiple benchmarks confirm the efficacy of LongRoPE. Notably, this framework preserves the original network design, requiring only minimal adjustments to the positional embedding scheme, and allows the reuse of established optimizations. Implementation will be accessible at \texttt{https://github.com/microsoft/LongRoPE}.

\section{Introduction}

Although Large Language Models (LLMs) have achieved impressive outcomes across a range of tasks~\cite{gpt4,llama2}, they are constrained by a finite context window size, such as LLaMA2’s maximum of 4096 tokens~\cite{llama2}. When inputs exceed this window, LLMs experience deteriorating performance due to unfamiliar positional embeddings not encountered during training. This limitation complicates key applications, including in-context learning with a substantial number of examples~\cite{Huang2023FewerIM} and deploying LLM agents~\cite{park2023generative,madaan2023selfrefine}.

Recent studies indicate that it is possible to enlarge the context window of a pre-trained LLM up to approximately 128k by conducting fine-tuning on longer passages~\cite{longlora,pi,yarn,zhang2024soaring,liu2023scaling}. However, three significant challenges hinder the further expansion of this window. Firstly, introducing new, untrained position indices often results in numerous problematic values, causing distributional shifts that make the fine-tuning process unstable and hinder convergence~\cite{pi}. This issue is particularly pronounced when scaling from 4k to over 1000k positions, as the majority (more than 90\%) of indices are untrained. Secondly, the fine-tuning procedure necessitates access to sufficiently long training data, but texts that surpass 1000k tokens are scarce in existing datasets. Additionally, training with such ultra-long sequences is highly resource-intensive, demanding extensive computation time and significant GPU investment. Thirdly, expanding to vast context windows exacerbates the dispersion of attention mechanisms, which must operate across an enormous number of token positions. This broad distribution diminishes performance in tasks involving shorter contexts, as previously reported~\cite{pi}.

A common method to address the initial challenge involves interpolating RoPE positional embeddings~\cite{rope,pi}, where newly assigned position indices are scaled down to fit within the pretrained boundaries, as depicted in Fig.\ref{fig:overview}. Position Interpolation (PI)~\cite{pi} implements a linear adjustment to RoPE’s rotary angle parameters, proportional to the expansion factor. NTK~\cite{ntk,dynamicntk} proposes using a non-linear scheme, combining both unequal interpolation and extrapolation over different dimensions in RoPE. YaRN~\cite{yarn} further refines this idea by segmenting RoPE dimensions based on frequency, subsequently employing extrapolation, NTK-style unequal interpolation, and standard linear interpolation across the respective groups. Nonetheless, the positional embeddings within Transformer networks demonstrate \emph{complex, highly variable information entropy} across dimensions. Current strategies fail to exploit this intricate variation fully, resulting in the omission of vital information and consequently restricting the achievable context window size.

Section~\ref{sec:analysis} empirically demonstrates two principal insights: \textbf{(1)} To perform effective positional interpolation, it is crucial to address two specific non-uniformities: the variations across RoPE dimensions and differences among token positions. Less interpolation proves advantageous for dimensions at the lower end of RoPE and for earlier token positions, although the optimal strategies are contingent on the desired extrapolation length.
\textbf{(2)} Incorporating these forms of non-uniformity into positional interpolation significantly aids in preserving essential information inherent to the original RoPE representation, especially with respect to critical dimensions and positional indices. This reduces the information loss ordinarily introduced by interpolation, thus offering improved initial states for subsequent fine-tuning. Additionally, this methodology supports extensions up to 8$\times$ in contexts where fine-tuning is not performed.

Building upon these insights, we present \textbf{LongRoPE}, a robust approach that pushes the LLM context window length past 2 \emph{million} tokens. The architecture of LongRoPE incorporates three principal contributions. To start, {LongRoPE} leverages complex multidimensional non-uniformities inherent in positional interpolation, enabling the determination of optimal rescale factors for RoPE’s rotation angles across each RoPE axis and adjusting according to token locations. Given that the search domain for these rescale factors expands exponentially with increasing extension ratios, {LongRoPE} deploys an evolutionary search methodology equipped with two optimization strategies to enhance the efficiency of this search. An illustration of the derived rescaled RoPE is provided in Fig.~\ref{fig:overview}.

Subsequently, {LongRoPE} utilizes an effective and incremental extension technique to reach a 2048k context window, circumventing the necessity for direct fine-tuning on extremely lengthy text data, which are both uncommon and difficult to obtain. The approach initiates by identifying a 256k context length within the pre-trained LLM, followed by fine-tuning at this length. Due to the non-uniform positional interpolation mechanism, which enables an 8$\times$ window expansion in settings where fine-tuning is not required, a secondary search is performed to determine optimal RoPE rescale factors for the newly fine-tuned extended LLM. Through this process, a 2048k context window is attained for models such as LLaMA2 and Mistral~\cite{mistral}.

In order to address potential drops in performance for the original, shorter context window, {LongRoPE} persistently tunes the RoPE rescale factors even after extending the LLM. Employing the same methodology as with expansion from 256k to 2048k, we reduce the context window size to 4k and 8k on the LLM fine-tuned at 256k, utilizing our search algorithm to decrease reliance on positional interpolation. When performing inference and the sequence length is below 8k, RoPE is updated according to the discovered rescale factors.

A comprehensive set of experiments involving multiple LLMs and a range of long-context tasks underscores the efficacy of our approach. The results reveal that LongRoPE consistently sustains low perplexity across evaluation lengths from 4k up to 2048k, attains a passkey retrieval accuracy exceeding 90\%, and demonstrates equivalent performance on standard benchmarks tailored to a 4096 context window. Furthermore, LongRoPE is compatible with any LLM utilizing RoPE embeddings. We intend to make our code and LongRoPE-2048k models publicly available.

\section{Non-uniformity in Positional Interpolation}
\label{sec:analysis}

\subsection{Preliminary}

Transformer models require explicit positional information, often in the form of position embedding, to represent the order of input tokens. Our work focuses on the RoPE~\cite{rope} position embedding, which is widely used in recent  LLMs. For a token at position index $n$, its corresponding RoPE encoding can be simplified as follows:

\begin{equation}
		\fontsize{7.8}{7.8} \selectfont
	\label{eq:rope}
	\left[ cos(n\theta_0), sin(n\theta_0), cos(n\theta_1), \cdot\cdot\cdot, cos(n\theta_{d/2-1}), sin(n\theta_{d/2-1}) \right]   
\end{equation}

where $d$ denotes the embedding size, 
$n\theta_i$ corresponds to the rotary angle assigned to the token in position $n$, 
and $\theta_i=\theta^{-2i/d}$ specifies the rotation frequencies. In the RoPE mechanism, the base value for $\theta$ is set to 10000 by default.

\textbf{\emph{Context window extension ratio $s$} and positional interpolation}. The ratio $s$ is defined as the proportion between the new (extended) context length $L'$ and the initial context length $L$: $s=\frac{L'}{L}$.

To extend the context window from $L$ to $L'$, current positional interpolation methods suggest downscaling rotation frequencies $\theta_i$ based on extension ratio $s$. Let $\beta=\theta^{2/d}$, and $\lambda$ denote the actual rescale factor related to $s$, we   unify these positional interpolation methods as follows:

\begin{equation}	
	\fontsize{7.5}{7.5} \selectfont
	\label{eq:unified_rope}
	\begin{aligned}
		\left[cos\left(\frac{n}{\lambda(\beta)^0}\right), sin \left(\frac{n}{\lambda(\beta)^0}\right), cos\left(\frac{n}{\lambda(\beta)^1}\right), \cdot\cdot\cdot,  sin \left(\frac{n}{\lambda(\beta)^{d/2-1}}\right) \right]   
	\end{aligned}
\end{equation}

\textbf{Linear positional interpolation (PI)}. PI~\cite{pi} proposes applying linear interpolation to position indices when operating within the original pre-training sequence length. When extending to a new target length by a ratio $s$, the rotation angles assigned to each position are scaled down uniformly by a factor $\lambda=s$ across all RoPE dimensions. Nonetheless, this approach compresses the positional information, causing nearby tokens to become less distinguishable from one another. Consequently, PI's effectiveness declines as the extension ratio increases.

\textbf{NTK-based interpolation and extrapolation}. ~\cite{ntk,dynamicntk} approach RoPE through the lens of information representation and leverage the framework of Neural Tangent Kernel (NTK) theory~\cite{jacot2018neural,tancik2020fourier}. To address the position saturation encountered in PI, they propose allocating the burden of interpolation across multiple RoPE dimensions. By applying less scaling to lower (high-frequency) dimensions and greater scaling to higher (low-frequency) ones, this approach supports both interpolation and extrapolation of positions, following the relationship $\lambda=s^{i}$. The dynamic NTK variant~\cite{dynamicntk} further refines this by adapting the extension factor at each position relative to the given sequence length. In contrast to PI, which relies on further training, NTK-informed techniques facilitate the extension of context windows without additional fine-tuning. However, these methods typically limit the feasible extension to at most 4$\times$.

\textbf{YaRN}~\cite{yarn} presents a notable advancement in positional interpolation by segmenting RoPE dimensions into three categories based on their frequency characteristics, and applying distinct interpolation methods to each. Specifically, extrapolation ($\lambda$=1) is utilized for the high frequency group, linear interpolation (PI) for the low frequency group, and the NTK approach for dimensions in the intermediate range. The essential innovation of YaRN is its frequency-based partitioning of RoPE dimensions; however, this grouping relies on manual empirical procedures, which could hinder its optimal effectiveness when adapting to novel LLM architectures.

\subsection{Study on Non-uniform Positional Interpolation}

Drawing motivation from NTK and YaRN, we observe improvements due to the introduction of non-linear functions—particularly when varying frequency treatments along RoPE dimensions that enable tailored interpolation and extrapolation. Despite this, existing methods predominantly depend on manually crafted heuristics. This observation leads to two key inquiries: (1) Can the current approach to positional interpolation be considered optimal? (2) Might there exist untapped non-linear strategies?

In order to address these inquiries, we employ evolution search (refer to Sec.~\ref{sec:method}) to identify improved non-uniform positional interpolations tailored for LLaMA2-7B. The optimization process is steered by perplexity metrics, utilizing 5 randomly selected instances from the PG19~\cite{pg19} validation dataset. Our experimental evaluation uncovers several important observations.

\textbf{Finding 1}: \textit{RoPE dimensions exhibit substantial non-uniformities, which are not effectively handled by current positional interpolation methods.}

We determine the optimal value of $\lambda$ for each RoPE dimension as described in Eq.~\ref{eq:unified_rope}. Table~\ref{tbl:exp1} presents a comparison of perplexity scores for LLaMA2-7B across different methods on the PG19 and Proof-pile~\cite{proof-pile} test datasets, all evaluated without fine-tuning. The results obtained from our optimized approach demonstrate substantial gains, highlighting the limitations of existing linear (PI) and non-uniform (Dynamic-NTK and YaRN) interpolation techniques. Interestingly, YaRN performs worse than PI and NTK on PG19 because it does not achieve the intended context window length when large language models are evaluated without fine-tuning. For instance, the perplexity associated with YaRN increases sharply beyond 7k tokens within an 8k context window.

Our investigation reveals that the rescaling coefficients $\lambda$ in Eq.~\ref{eq:unified_rope} are heterogeneous, in contrast to the constant factor $s$ adopted in PI, the formulaic approach of NTK, or the group-level computation used in YaRN. Employing these variable scaling factors leads to notable enhancements in LLaMA2’s language modeling accuracy (perplexity) over 8k and 16k context lengths, all without additional fine-tuning. This improvement arises because the resulting positional encodings retain key aspects of the original RoPE structure—most crucially in certain dimensions—thus facilitating the model's ability to discern adjacent token positions with greater ease.

\textbf{Finding 2}: \textit{RoPE for the initial tokens in the input sequence should be extrapolated with less interpolation.} 

For the first $\hat{n}$ tokens within the input sequences, we propose that their RoPE should undergo minimal interpolation. The rationale is that these tokens tend to receive the highest attention scores and play a pivotal role in attention layers, as demonstrated in prior work such as Streaming LLM~\cite{streamingllm} and LM-Infinite~\cite{han2023lminfinite}. To examine this hypothesis, we expanded the context window to 8k and 16k using both PI and NTK approaches, while deliberately preserving the first $\hat n$ tokens (with $\hat n$ ranging from 0, 2, ..., 256) without interpolation. Setting $\hat n$=0 results in standard PI and NTK. Table~\ref{tbl:tokenposition} reveals two principal findings: \textbf{(1)} Excluding the initial tokens from position interpolation leads to measurable performance gains. \textbf{(2)} The most effective value for $\hat n$, or the number of early tokens left un-interpolated, is contingent on the desired extension length.

\textbf{Finding 3}: \textit{Non-uniform positional interpolation effectively extends LLM context window in both fine-tuning and non-fine-tuning settings.} 

Although our explored non-uniform position interpolation technique demonstrates notable improvements in extension performance at 8k and 16k context lengths without requiring fine-tuning, achieving satisfactory results for even longer contexts necessitates a fine-tuning phase. Therefore, we conduct fine-tuning of LLaMA2-7B using the RoPE configuration derived from our search for a 64k context window (refer to Appendix for detailed configurations). As presented in Table~\ref{tbl:finetune}, our approach consistently exceeds the performance of PI and YaRN, both prior to and following fine-tuning of LLaMA2-7B. This enhanced performance is attributed to our approach’s effective application of non-uniform positional interpolation, which better preserves information and offers a more optimal starting point for subsequent fine-tuning.

\textbf{Summary}. This work identifies two forms of non-uniformity—across RoPE dimensions and among token positions. Appropriately leveraging these non-uniformities in the interpolation of positions yields substantial gains in the ability of LLMs to extend their context windows.

\section{LongRoPE}

\label{sec:method}
Inspired by these results, we propose LongRoPE, which incorporates a novel, efficient search procedure to comprehensively harness the identified non-uniformities, subsequently leveraging this method to expand the context window for LLMs past 2 million tokens.

\subsection{Problem Formulation}

These two sources of non-uniformity significantly expand the solution space and add layers of difficulty to the optimization process. In response, we conceptualize the optimization challenge of multidimensional non-uniform position interpolation as a search problem.

For a LLM targeting a  context window size of $L'$ and lengthy input documents $\mathbf{X}$, where each $\mathbf{x}\in\mathbf{X}$ surpasses $L'$ in token length, we denote the original rotary angle of the $i^{th}$ dimension in RoPE embedding at token position $n$ as  $\frac{n}{\beta^i}$. The optimization problem is then formulated as follows:

\begin{equation}
\begin{gathered}
		\label{eq:problem}

\underset {\mathbf{x}\in\mathbf{X}; \, |\mathbf{x}|\ge L'}{ \operatorname {arg\,min} } \,  \mathcal{L}\, \left( \text{LLM} (\text{RoPE}, \mathbf{X})\right),	\text{where } 
\\
\scalebox{0.93}{
$\underset {\begin{subarray}{l} i=0,\cdot\cdot,\frac{d}{2}-1;\\ n\in[ 0,|\mathbf{x}|);\end{subarray}}{\text{RoPE}_(n)}= {\left[\cdots, \cos\left(\mathbb{I}(\hat{\lambda}_{i}, \hat{n})\cdot\frac{n}{\beta^i}\right),  \sin\left(\mathbb{I}(\hat{\lambda}_{i}, \hat{n})\cdot\frac{n}{\beta^i}\right), \cdots \right] }$}\\
	\text{where } \mathbb{I}(\hat{\lambda}_{i},\hat{n})=\begin{cases}
	1&\mbox{ for $n<\hat{n}$}\\
{\frac{1}{\lambda}_{i}}&\mbox{ if $n\geq\hat{n}$}
\end{cases}
	\end{gathered}
\end{equation}

In this framework, we implement a set of rescaling coefficients, $\mathbb{I}(\hat{\lambda}_{i},\hat{n})$, designed to address two distinct types of non-uniformities. The variables $\hat{\lambda_i}$ and $\hat{n}$ correspond, respectively, to irregularities found in the RoPE dimension and in token position. In detail, $\mathbb{I}(\hat{\lambda}_{i},\hat{n})$ adjusts the rotation angle specifically for the $i^{th}$ RoPE dimension; here, $\hat\lambda_{i}$ acts as the scaling parameter, while $\hat{n}$ serves as the threshold for token position. For token positions earlier than $\hat{n}$ (specifically, the first $\hat{n}$-1 tokens), the original RoPE rotary angle $\frac{n}{\beta^i}$ remains unmodified by $\hat\lambda_{i}$. When the token positions satisfy $n \ge \hat{n}$, the rescaling factor is then applied.

With a specified target context window length of $L'$, the task is to determine the most effective rescale factors ($\mathbb{I}(\hat{\lambda}_{0},\hat{n})$, $\mathbb{I}(\hat{\lambda}_{1},\hat{n})$, ..., $\mathbb{I}(\hat{\lambda}_{i},\hat{n})$, ...) for each of the $d$ RoPE dimensions (from the $1^{st}$ up to the $d^{th}$). By applying these optimized rescale factors to the $\text{RoPE}$ mechanism within the $\text{LLM}$, the goal is to minimize the next-token prediction loss $\mathcal{L}$—equivalently, the perplexity—when processing input sequences $\mathbf{X}$ of token length $L'$.

\subsection{Searching the Non-uniform Position Interpolation}

To address the challenge posed in Eq.~\ref{eq:problem}, we present a straightforward but powerful approach that identifies the best RoPE rescale factors, allowing for comprehensive utilization of the multidimensional variations inherent in position embedding.

\textbf{Search space}. We construct an extensive search space that encompasses both forms of non-uniformity. The configuration of this space is depicted in Table~\ref{tbl:searchspace}. More precisely, our approach enables the optimization of an individual rescale factor for each dimension within the RoPE embedding. To streamline the design of the search space, we opt to explore $\lambda_i$ and $\hat{n}$ instead of directly optimizing $\mathbb{I}(\hat{\lambda}_{i},\hat{n})$, where  $\hat{\lambda}_{i}=1/\lambda_i$. As presented in Table~\ref{tbl:searchspace}, the permissible range for $\lambda_i$ begins at 1.0 (which corresponds to basic extrapolation) and extends up to $s\times1.25$ (representing more pronounced interpolation than that used in PI), with increments of 0.01. Here, $s$ denotes the desired ratio for context window enlargement.

The parameter $\hat{n}$ determines how many initial token positions preserve their original positional encoding, foregoing any position interpolation (specifically, maintaining the native RoPE embedding). In practice, $\hat{n}$ is chosen from the set \{0, 1, 2, 4, 8, 12, 16, 20, 24, 28, 32, 64, 128, 256\}. Setting $\hat{n}=0$ indicates that every token position applies the identified rescale factors.

\textbf{Evolution-based search}. The search space described in Table~\ref{tbl:searchspace} encompasses a wide array of possible positional interpolation configurations, making thorough exploration highly demanding. For instance, when considering an extension factor of $s=4\times$, this results in $400^{128/2}\times14$=4$\times10^{167}$ potential options. As the extension ratio increases, the search space grows at an exponential rate. To effectively navigate this complexity, we employ evolution search~\cite{guo2020single} and introduce two optimization strategies that significantly enhance the efficiency of the search process. The steps of the overall search are depicted in Algorithm~\ref{alg:search}.

\textit{Optimized initial population generation}. In place of fully random initialization of a population containing $P$ rescale factors, our approach explicitly includes the RoPE rescale factors for PI, NTK, and YaRN among the initial set of individuals. The rest of the population, consisting of $P$-3 members, is produced by randomly introducing mutations to these three rescale factors with probability $p$.

\textit{Monotonically non-decreasing constraint}. Following the creation of the initial population, each candidate is evaluated using LLM perplexity. This involves adjusting the target LLM with the associated RoPE rescale factors and then determining the perplexity for the input $\mathbf{X}$. The highest-performing k candidates are selected as parents for the next evolutionary step. Due to the immense breadth of the search space, simplistic mutation and crossover may produce suboptimal solutions, resulting in excessive and often redundant perplexity evaluations. This inefficiency is exacerbated as $L'$ grows larger, since each perplexity measurement requires costly inference.

To mitigate this issue, we enforce a monotonicity constraint such that the RoPE rescaled factors sampled are non-decreasing: $\lambda_i\le \lambda_{i+1}$. Only those RoPE parameterizations satisfying this property are considered for perplexity testing within the LLM, which notably reduces the computational expense of searching. Concretely, we stipulate that $\lambda_i$ grows monotonically along with the RoPE dimension indices ($i = 0, ..., 63$). This choice of monotonicity by dimension is informed by NTK theory~\cite{jacot2018neural,tancik2020fourier,ntk}, which indicates that lower-dimensional, high-frequency components necessitate less interpolation (implying a smaller $\lambda_i$), whereas higher-dimensional, lower-frequency components can support more interpolation (thus allowing for larger $\lambda_i$).

\begin{algorithm}[t]
	\small
	\caption{The search algorithm}
	\textbf{Input:} target LLM, input samples $\mathbf{X}$, population size $P$, mutation size $N_1$, crossover size $N_2$, maximum iterations $\mathcal{T}$, mutation probability $p$\\

	\begin{algorithmic}[1]
		\label{alg:search}
		\STATE Top-k $\gets$ $\phi$;	
		\STATE $\text{P}_0 \gets$ \textit{Initialize\_population\_with\_optimization}($P$, $\mathbf{X}$, $p$);
		\FOR{$i=1$ to $\mathcal{T}$}
		\STATE \textit{Compute\_perplexity}(LLM, $\text{P}_{i-1}$, $\mathbf{X}$);
		\STATE Top-k $\gets$ \textit{Update\_Topk}(Top-k, $\text{P}_{i-1}$);
		\STATE $\text{P}_{mutation} \gets$ \textit{Mutation\_with\_mono\_constraint}(Top-k, $N_1$, $p$);
		\STATE $\text{P}_{crossover} \gets$ \textit{Crossover\_with\_mono\_constraint}(Top-k, $N_2$);
		\STATE $\text{P}_i \gets \text{P}_{mutation} \cup \text{P}_{crossover} \cup$ Top-k;
		\ENDFOR
		\STATE Output the member in Top-k exhibiting the minimum perplexity;
	\end{algorithmic}
\end{algorithm}

\textbf{8$\times$ extension without fine-tuning}. Through our evolutionary search approach, we discover optimal, non-uniform RoPE rescale factors that strategically maintain essential dimensions and positional information, thereby reducing the loss incurred by interpolation. As illustrated in Fig.\ref{fig:non-ft}, this technique successfully scales the context window of LLaMA2 from 4k to 32k tokens absent any fine-tuning. Conversely, approaches like PI, as well as non-uniform variants of NTK and YaRN, experience a sharp increase in perplexity beyond a 2$\times$ window extension.

\subsection{Extending LLM Context Window to 2048K}
\label{sec:extendto2048k}

\textbf{Progressive extension to 2048k}. In this section, we present our approach for scaling the context window of pre-trained LLMs from the standard 4k tokens up to beyond 2048k tokens. Our findings show that non-uniform positional interpolation is sufficient to obtain an 8$\times$ increase in window size without necessitating additional fine-tuning. When seeking even greater extensions (e.g., by a factor of 512$\times$), however, fine-tuning becomes indispensable. A commonly adopted strategy is to identify appropriate RoPE rescaling parameters for the intended 2048k window, followed by a fine-tuning stage. Nevertheless, this approach is hampered by the extremely high computational demands required. Furthermore, our experiments indicate that achieving successful fine-tuning with such large extension ratios poses considerable difficulties (refer to the Appendix for further details).

LongRoPE demonstrates effectiveness when applied to both base and fine-tuned expanded LLMs. Accordingly, we propose a streamlined, incremental approach that attains the desired 2048k context window using only 1k fine-tuning iterations, all conducted with training sequences limited to 256k tokens.

$\Diamond$ \textit{Exploring the extension of pre-trained LLMs to 256k using LongRoPE search}. Using LLaMA2 as a case study, we perform searches for desired context window lengths of 128k and 256k. At this phase, the respective extension factors are 32$\times$ and 64$\times$.

$\Diamond$ \textit{Fine-tuning to 256k}. In this stage, the pre-trained LLM undergoes fine-tuning to reach a context window size of 256k. The process begins by fine-tuning LLaMA2 with RoPE rescaled factors set for 128k, completing 400 steps. Next, the RoPE rescaled factors are switched to 256k, and a subsequent 600 fine-tuning steps are performed on the existing checkpoint. This two-stage approach has been found to be more efficient compared to initializing fine-tuning directly at 256k.

$\Diamond$ \textit{Expanding fine-tuned extended LLMs to 2048k via LongRoPE search}. In the final phase, we conduct an additional search on the LLM that has already been fine-tuned to a context length of 256k. This approach yields a remarkably large context window of 2048k, achieved without any need for subsequent fine-tuning. Consequently, the resulting extension factor is $512\times$.

\textbf{Restoring performance in shorter context windows}. Upon expanding the context window to a massive 2048k tokens, a decline in performance is observed within the initial context window range. This phenomenon has been documented in positional interpolation literature~\cite{pi}, where adapting position embeddings to higher-dimensional spaces within the original context constrains them to a limited subspace, impairing model efficacy. When employing a 512$\times$ extension ratio, token positions inside the original 4k window are compressed significantly, leading to dense packing and diminished effectiveness.

To address this issue, we conduct an additional evolutionary search on the expanded LLM to fine-tune RoPE rescale factors tailored to short context lengths (such as 4k and 8k). The upper limit for searched $\lambda$ is decreased, as shorter sequences require less positional interpolation. When running inference, the LLM modifies the appropriate RoPE rescale factors in real time.

\section{LongRoPE}

\label{sec:method}
Building upon these observations, we propose LongRoPE. This approach initially devises an effective search method that leverages both types of non-uniformities, and subsequently applies this technique to push the LLM context window to exceed 2 million tokens.

\subsection{Problem Formulation}

These two forms of non-uniformity significantly expand the possible solution space and complicate the optimization process. To mitigate this issue, we recast the optimization challenge of multidimensional non-uniform position interpolation into a search-based problem.

For a LLM targeting a  context window size of $L'$ and lengthy input documents $\mathbf{X}$, where each $\mathbf{x}\in\mathbf{X}$ surpasses $L'$ in token length, we denote the original rotary angle of the $i^{th}$ dimension in RoPE embedding at token position $n$ as  $\frac{n}{\beta^i}$. The optimization problem is then formulated as follows:

\begin{equation}
\begin{gathered}
		\label{eq:problem}

\underset {\mathbf{x}\in\mathbf{X}; \, |\mathbf{x}|\ge L'}{ \operatorname {arg\,min} } \,  \mathcal{L}\, \left( \text{LLM} (\text{RoPE}, \mathbf{X})\right), \text{such that} 
\\
\scalebox{0.93}{
$\underset {\begin{subarray}{l} i=0,\ldots,\frac{d}{2}-1;\\ n\in[ 0,|\mathbf{x}|);\end{subarray}}{\text{RoPE}_{(n)}}= {\left[\ldots,cos\left(\mathbb{I}(\hat{\lambda}_{i}, \hat{n})\frac{n}{\beta^i}\right),  sin\left(\mathbb{I}(\hat{\lambda}_{i}, \hat{n})\frac{n}{\beta^i}\right),\ldots\right] }$}
\\
	\text{where} \quad \mathbb{I}(\hat{\lambda}_{i},\hat{n})=
	\begin{cases}
	1 & \text{if}~ n<\hat{n}\\
	\frac{1}{\lambda_i} & \text{if}~ n\geq\hat{n}
	\end{cases}
	\end{gathered}
\end{equation}

In this approach, we define a group of scaling factors, $\mathbb{I}(\hat{\lambda}_{i},\hat{n})$, to address the two distinct types of non-uniformity. The terms $\hat{\lambda_i}$ and $\hat{n}$ refer to irregularities found in RoPE dimensions and token positions, respectively. The scaling function $\mathbb{I}(\hat{\lambda}_{i},\hat{n})$ modifies the rotation angle associated with the $i^{th}$ RoPE dimension; $\hat\lambda_{i}$ acts as the multiplier while $\hat{n}$ specifies the position threshold for tokens. For token positions from $0$ up to $\hat{n}-1$, the scaling factor $\hat\lambda_{i}$ is omitted, and the conventional RoPE rotary angle $\frac{n}{\beta^i}$ remains unchanged. Once token positions reach $n\ge\hat{n}$, the scaling factor comes into effect.

With a desired context window length of $L'$, our goal is to determine the set of optimal rescale factors ($\mathbb{I}(\hat{\lambda}_{0},\hat{n})$, $\mathbb{I}(\hat{\lambda}_{1},\hat{n})$, ... , $\mathbb{I}(\hat{\lambda}_{i},\hat{n})$, ...) across the $1^{st}$ through $d^{th}$ dimensions of the RoPE mechanism. In doing so, we aim for the adjusted target $\text{LLM}$, employing the modified $\text{RoPE}$, to yield the lowest possible next token prediction loss, $\mathcal{L}$ (that is, perplexity) for input sequences $\mathbf{X}$ of length $L'$.

\subsection{Searching the Non-uniform Position Interpolation}

To address the challenge described in Eq.~\ref{eq:problem}, we propose a straightforward and robust approach that identifies ideal RoPE rescale factors, thereby taking full advantage of the complex, multidimensional variability present in position embedding.

\textbf{Search space}. We establish an extensive search space encompassing both forms of non-uniformity. The layout of this search space is presented in Table~\ref{tbl:searchspace}. In particular, we permit the optimization of distinct rescaling factors for every dimension within the RoPE embedding. To streamline the structure of the search space, rather than directly optimizing $\mathbb{I}(\hat{\lambda}_{i},\hat{n})$, we instead focus on $\lambda_i$ and $\hat{n}$, recognizing that $\hat{\lambda}_{i}=1/\lambda_i$. According to Table~\ref{tbl:searchspace}, $\lambda_i$ can be selected from 1.0 (representing pure extrapolation) up to $s\times1.25$ (representing a broader interpolation range than PI), in increments of 0.01. Here, $s$ signifies the expansion ratio for the context window being targeted.

The parameter $\hat{n}$ determines how many leading token positions keep their original RoPE embedding without any position interpolation. In our experiments, we let $\hat{n}$ vary over the set \{0, 1, 2, 4, 8, 12, 16, 20, 24, 28, 32, 64, 128, 256\}. If $\hat{n}=0$, every token position utilizes the optimized rescale factors.

\textbf{Evolution-based search}. The search space outlined in Table~\ref{tbl:searchspace} encompasses a vast array of positional interpolation options, making exhaustive exploration highly complex. For instance, when applying an extension ratio of $s=4\times$, the resulting configuration space becomes $400^{128/2}\times14$=4$\times10^{167}$ alternatives. As the extension ratio increases, the dimensionality of the search space grows at an exponential rate. To efficiently navigate this extensive landscape, we adopt evolution search~\cite{guo2020single} and implement two optimization strategies that significantly enhance the speed of the search process. The complete search workflow is depicted in Algorithm~\ref{alg:search}.

\textit{Optimized initial population generation}. Rather than filling the initial population of $P$ rescale factors entirely at random, we explicitly include the three RoPE rescale factors representing PI, NTK, and YaRN among the initial individuals. The other $P$-3 individuals are produced by introducing random mutations to these three rescale factors, each with probability $p$.

\textit{Monotonically non-decreasing constraint}. Once the initial population has been formed, each individual's perplexity is measured using the target LLM. To do so, we assign the specific RoPE rescale factors associated with each candidate to the model and evaluate perplexity over input $\mathbf{X}$. The $k$ individuals with the lowest perplexity are selected as parents for the next evolutionary steps. Due to the considerable breadth of the search space, basic mutation and crossover can frequently produce suboptimal candidates, resulting in redundant perplexity evaluations. This inefficiency becomes more pronounced for large values of $L'$, as each perplexity evaluation is computationally intensive.

To tackle this issue, we enforce a monotonicity constraint such that the rescaled RoPE factors adhere to $\lambda_i \le \lambda_{i+1}$, ensuring they do not decrease. Only RoPE configurations meeting this requirement are used for LLM perplexity evaluation, thereby greatly streamlining the search process. Concretely, we mandate that $\lambda_i$ rises monotonically as the RoPE dimension (with $i$ ranging from 0 to 63) increases. This imposition of monotonicity across dimensions draws upon NTK theory~\cite{jacot2018neural,tancik2020fourier,ntk}, which argues that lower-indexed dimensions associated with higher frequencies benefit from minimal interpolation (represented by a smaller $\lambda_i$), while higher-indexed dimensions with lower frequencies accommodate more extensive interpolation (reflected by larger $\lambda_i$ values).

\begin{algorithm}[t]
	\small
	\caption{The search algorithm}
	\textbf{Input:} target LLM, input samples $\mathbf{X}$, population size $P$, mutation size $N_1$, crossover size $N_2$, maximum iterations $\mathcal{T}$, mutation probability $p$\\

	\begin{algorithmic}[1]
		\label{alg:search}
		\STATE Top-k $\leftarrow$ $\phi$;	
		\STATE $\text{P}_0$ $\leftarrow$ \textit{Initialize\_population\_with\_optimization}($P$, $\mathbf{X}$, $p$); 
		\FOR{$i$ = $1$ \textbf{to} $\mathcal{T}$}
		\STATE \textit{Compute\_perplexity} (LLM, $\text{P}_{i-1}$, $\mathbf{X}$);
		\STATE Top-k $\leftarrow$ \textit{Update\_Topk}(Top-k, $\text{P}_{i-1}$);
		\STATE $\text{P}_{mutation}$ $\leftarrow$ \textit{Mutation\_with\_mono\_constraint}(Top-k, $N_1$, $p$); 
		\STATE $\text{P}_{crossover}$ $\leftarrow$ \textit{Crossover\_with\_mono\_constraint}(Top-k, $N_2$);
		\STATE $\text{P}_i$ $\leftarrow$ $\text{P}_{mutation}$ $\cup$ $\text{P}_{crossover}$ $\cup$ Top-k;
		\ENDFOR
		\STATE Output the member with the minimal perplexity from Top-k;
	\end{algorithmic}
\end{algorithm}

\textbf{8$\times$ extension without fine-tuning}. Through evolutionary search, our approach efficiently discovers optimal non-uniform RoPE rescaling parameters, strategically retaining critical dimensions and positions to reduce information loss stemming from interpolation. As illustrated in Fig.\ref{fig:non-ft}, this technique succeeds in expanding LLaMA2’s context window from 4k up to 32k tokens without requiring any fine-tuning. By comparison, methods like PI, as well as non-uniform NTK and YaRN, exhibit substantial increases in perplexity once the extension exceeds 2$\times$.

\subsection{Extending LLM Context Window to 2048K}
\label{sec:extendto2048k}

\textbf{Progressive extension to 2048k}. Here, we present our approach for expanding the context window of pre-trained LLMs from the conventional 4k to beyond 2048k. Our experiments show that using non-uniform positional interpolation enables up to an 8$\times$ context window increase without the need for fine-tuning. When more significant extensions are desired, such as 512$\times$, fine-tuning becomes essential. A typical strategy involves identifying appropriate RoPE rescaling parameters at the target 2048k length, followed by fine-tuning. This approach, however, is often impractical because of the substantial computational resources required. Additionally, our experience indicates that effective fine-tuning with such high extension ratios is quite difficult (see Appendix for details).

Fortunately, LongRoPE proves to be efficient for both pre-trained and fine-tuned extended LLMs. As a result, we propose a streamlined, incremental approach capable of reaching the desired 2048k target through only 1k fine-tuning iterations, all within a training sequence length limit of 256k.

$\Diamond$ \textit{Expanding pre-trained LLM to 256k using LongRoPE search}. Using LLaMA2 for illustration, we perform a search procedure for desired context window lengths of 128k and 256k tokens. At these levels, the context window is increased by factors of 32$\times$ and 64$\times$, accordingly.

$\Diamond$ \textit{Fine-tuning to 256k}. Subsequently, we adapt the pre-trained LLM for a 256k context window through a multi-stage fine-tuning process. Initially, we fine-tune LLaMA2 for 400 steps utilizing the RoPE rescaled factors adjusted for 128k. Following this, we substitute the RoPE rescaled factors with those suitable for 256k in the obtained checkpoint, and continue with an extra 600 fine-tuning steps. This approach demonstrates greater efficiency compared to immediate fine-tuning for 256k.

$\Diamond$ \textit{LongRoPE search enables increasing the context length of a fine-tuned extended LLM up to 2048k}. In this step, we conduct a subsequent search process on the fine-tuned model originally trained for a 256k-length context window. This approach achieves an exceptionally large context window size of 2048k, and notably, does not necessitate additional rounds of fine-tuning. The overall enlargement factor attained is $512\times$.

\textbf{Restoring performance in shorter context windows.} When scaling up to a very large 2048k context window, we observe a degradation in model effectiveness within the range of the initial context window. This phenomenon has been previously identified in studies on positional interpolation~\cite{pi}, where the compression of position embeddings into a much tighter segment within the original window leads to diminished model performance. The issue becomes pronounced with a 512$\times$ scaling factor, as positions mapped inside the initial 4k context are densely packed.

To address this issue, an additional evolution search is conducted on the extended LLM to fine-tune the RoPE rescale factors specifically for shorter context lengths, such as 4k and 8k. Since less positional interpolation is necessary at these reduced lengths, the upper bound for the searched $\lambda$ is lowered accordingly. When performing inference, the LLM automatically adapts the relevant RoPE rescale factors.

 \section{Experiments}

\label{sec:exp}
\subsection{Setup}
\textbf{Evaluation Tasks and models}. LongRoPE is utilized with both LLaMA2-7B and Mistral-7B models. Our evaluation encompasses three primary dimensions: (1) the perplexity exhibited by extended-context LLMs when processing lengthy documents; (2) a Passkey retrieval task designed to assess how well the model extracts a specific passkey embedded in abundant irrelevant content; and (3) performance on established LLM benchmarks restricted to a 4096-token context window.

\textbf{Fine-tuning}. LLaMA2 is fine-tuned using a learning rate of $2\times10^{-5}$, which is linearly decayed, and a global batch size set at 32. The model undergoes 400 training steps utilizing the Redpajama~\cite{redpajama} dataset, which is divided into 128k-length segments, each framed with BOS and EOS tokens. After reaching an initial checkpoint, training continues for another 600 steps to extend the context window to 256k tokens. Fine-tuning at 128k context is conducted across 8 A100 GPUs via the distributed training system~\cite{cube}, whereas extending to 256k tokens necessitates 16 A100 GPUs. For Mistral, training employs a fixed learning rate of $1\times10^{-6}$ alongside a global batch size of 64. For both the 128k and 256k variants, the procedure follows YaRN~\cite{yarn} settings: 400 steps are performed on Together Computer's Long-Data Collections~\cite{mistraldata}, with a sequence length of 16k. Model training is carried out using 4 A100 GPUs.

\textbf{Search}. For target window sizes up to 256k, we set $P$=64, $N_1$=$N_2$=16, $p$=0.3, and $\mathcal{T}$=40, while selecting the top-32 candidates for mutation and crossover in each generation. Perplexity assessments are conducted using 5 randomly chosen samples from the PG19 validation set, ensuring that each sample is at least as long as the desired context window. For window sizes greater than 512k, both the population and the numbers of mutations and crossovers per iteration are reduced by half. In this regime, perplexity is determined using 3 random samples from the Pile-Books3~\cite{pile} validation set.

\textbf{Baselines}. In order to achieve a context window size of 2048k, we performed fine-tuning on models with initial context lengths of 128k and 256k. This results in LongRoPE-2048k (ft=128k) and LongRoPE-2048k (ft=256k) corresponding to LLaMA2 and Mistral, respectively. We assess these four models against leading approaches for enlarging context windows, focusing on publicly available LLMs that have undergone additional fine-tuning following positional interpolation via PI, NTK, and YaRN techniques. The evaluated baselines encompass Together-32k~\cite{together}, Code LLaMA~\cite{codellama}, LongLoRA-full-FT-100k~\cite{longlora}, as well as YaRN-LLaMA and YaRN-Mistral~\cite{yarn}.

\subsection{Main Results}

\label{sec:mainresult}
\textbf{Language modeling over extended contexts up to 256k.} Our analysis commences with a comparison to leading long-context LLMs for evaluation sequences reaching 256k tokens. To illustrate the broad applicability of our approach, we employ two different datasets: the test splits of Proof-pile~\cite{pg19} and PG19~\cite{pile}. Perplexity is assessed across multiple context sizes through a sliding window method with a window size of 256 tokens. For PG19, evaluation is conducted over the full test set containing 100 documents. In the case of Proof-pile, we adopt the protocol described in YaRN~\cite{yarn}, randomly selecting 10 samples, each with a minimum sequence length of 128k.

Table~\ref{tbl:proof-pile} and Table~\ref{tbl:pg19} present a comparison of perplexity scores for LLaMA2 and Mistral models that have been augmented using various interpolation techniques on the Proof-pile and PG19 benchmarks, respectively. Two principal findings emerge: \textbf{(1)} As the evaluation length increases from 4k up to 256k, the perplexity of our extended models consistently decreases, demonstrating their enhanced capacity to utilize extended context. \textbf{(2)} Despite extending the context window by a factor of 16—a scenario that usually poses difficulties for sustaining performance on shorter sequences—our LongRoPE-2048k models still achieve better results than the latest state-of-the-art baselines at the 256k context length.

\textbf{Long sequence language modeling beyond 2000k}. To assess model performance on very large-scale texts, we utilize the Books3~\cite{pile} dataset. To ensure evaluation remains efficient, 20 books with lengths greater than 2048k are sampled at random, and a sliding window approach with a size of 256k is applied.

As detailed in Table~\ref{tbl:books3}, LongRoPE effectively expands the context window for both LLaMA2-7B and Mistral-7B models to 2048k, while maintaining or surpassing the baseline perplexity values for shorter contexts ranging from 8k to 128k. Significant disparities are apparent between LLaMA2 and Mistral at the 2048k context length. While Mistral demonstrates superior performance at reduced window sizes, its perplexity rises above 7 for lengths greater than 256k. LLaMA2’s behavior is consistent with anticipated outcomes: its perplexity steadily declines as the window length increases, although it experiences slight increases at the 1024k and 2048k marks. Additionally, LongRoPE-2048k yields improved results when fine-tuned at 256k for LLaMA2 compared to fine-tuning at 128k, attributable to the lower secondary extension ratio (8$\times$ versus 16$\times$). Conversely, Mistral achieves more favorable outcomes when fine-tuned at 128k. This phenomenon primarily stems from the use of a 16k training length, as prescribed by YaRN's methodology, during the 128k and 256k fine-tuning of Mistral, which limits the model’s capacity for subsequent extension of the context window.

\textbf{Passkey retrieval}. Here, we investigate the practical limits of context window size for generation tasks by utilizing a synthetic benchmark introduced by ~\cite{passkey}. Specifically, the task requires the model to locate and reproduce a randomly selected passkey—represented as a five-digit number—that is embedded somewhere within an extended document. The specifics of the prompt format are provided in the appendix. For evaluation, we conduct 10 repetitions of the passkey retrieval scenario, ensuring that the passkey appears at uniformly random positions throughout the entire context window used for assessment.

As depicted in Fig.~\ref{fig:passkey}, retrieval accuracy is compared across various baseline models. The performance of existing approaches declines sharply, reaching zero accuracy past 128k tokens. In contrast, our LongRoPE-LLaMA2-2048k (ft=256k) achieves consistently high retrieval accuracy ($\ge$90\%) throughout the range from 4k to 2048k, even under the demanding condition of retrieving a passkey among millions of tokens. Furthermore, LongRoPE-Mistral-2048k (ft=128k) sustains perfect accuracy up to 1800k tokens, only decreasing to 60\% at 2048k. This result is consistent with Table~\ref{tbl:books3}, which shows a slight rise in perplexity at the 2048k token length.

\textbf{Performance assessment on standard benchmarks in the native context window}. For evaluating LongRoPE-2048k models within their standard context length, we utilize the Hugging Face Open LLM Leaderboard~\cite{huggingfaceboard} to measure outcomes across both zero-shot and few-shot scenarios. Specifically, we adopt a 25-shot setting for ARC-Challenge~\cite{allenai:arc}, 10-shot evaluation on HellaSwag~\cite{zellers2019hellaswag}, a 5-shot approach for MMLU~\cite{mmlu}, and a zero-shot configuration for TruthfulQA~\cite{lin2021truthfulqa}.

Table~\ref{tbl:commonsense} illustrates that our models deliver similar performance to baselines on the initial benchmark, which targets shorter context windows. Notably, our models even surpass the original Mistral on TruthfulQA by a margin of +0.5\%. While LongRoPE-LLaMA2-2048k, fine-tuned at a context length of 256k, exhibits a modest increase in performance drop, its results generally remain acceptable across the evaluated tasks.

\subsection{Ablation Results}

\textbf{Assessing the impact of the second positional interpolation}. Within the progressive extension framework, we employ our search algorithm to perform an additional, non-uniform positional interpolation on LLMs that have already undergone fine-tuning and extension. We assess the efficacy of this step using our fine-tuned LLaMA2-256k model, which we further extend to 512k, 1024k, and 2048k sequence lengths using both PI and YaRN. Results in Table~\ref{tbl:secondsearch} demonstrate that our approach to non-uniform positional interpolation preserves a stable perplexity level throughout the extension process. In comparison, perplexity for PI and YaRN rises sharply as the extension factor increases.

\textbf{Recovery efficiency for reduced context lengths.} In order to counteract the decline in performance at smaller context lengths, we fine-tune the RoPE parameters for LongRoPE-2048k using our search methodology. The adjustment entails lowering the upper limit for scaling factors during the search process, thus favoring reduced interpolation when evaluating 4k and 8k contexts. Table~\ref{tbl:shortperformance} presents the perplexity metrics for LongRoPE-LLaMA2-2048k applied to Proof-pile datasets at 4k and 8k context sizes, together with mean scores from LLM benchmarks. These findings strongly indicate marked gains in performance at the shorter context ranges.

\textbf{Examination of the two types of non-uniformity}. To assess the individual impact of each form of non-uniformity on model outcomes, we conducted ablation studies. Specifically, two experimental setups were utilized: (i) we extended LLaMA2-7B to shorter sequence lengths of 16k and 32k using three approaches—PI, searching solely for the optimal RoPE dimension, and simultaneously optimizing both non-uniformities; (ii) we expanded our fine-tuned LLaMA2 model from 256k to 2048k tokens, maintaining a consistent methodology. Perplexity measurements were taken without performing additional fine-tuning. As illustrated in Table~\ref{tbl:searchdimension}, introducing non-uniformity within the RoPE dimension leads to a notable reduction in perplexity relative to the linear interpolation employed by PI. Moreover, the incorporation of non-uniformity in token position enhances results at 16k and 32k, though its efficacy diminishes at the extreme length of 2048k—likely attributable to the sheer sequence length. The strategy of retaining only the initial tokens without interpolation does not prove effective in this context, and we suggest that this direction merits further investigation.

\section{Related Works}

Besides position interpolation methods, this section reviews alternative techniques explored in recent literature.

\textbf{Retrieval-based} methods leverage external memory components to store distant contextual information, utilizing retrieval mechanisms during inference to access relevant documents~\cite{longllama,wang2023augmenting,borgeaud2022improving}. Such strategies usually require deliberate architectural changes to the LLM itself. By comparison, our approach employs only minimal adjustments to positional embeddings, making it notably more lightweight. Furthermore, we support a broader array of long context applications beyond document retrieval, including tasks like extended summarization and few-shot learning.

\textbf{Attention-based context window extensions}. In addition to leveraging positional embedding interpolation, other works have managed to extend the input context by retaining the standard LLM context window length and instead adjusting the attention mechanism~\cite{han2023lminfinite, streamingllm,parallelwindow}. The central approach involves alleviating the problem of attention computation scaling arising from novel positions through innovative attention mask designs. Importantly, these strategies serve as a complement to positional interpolation techniques.

\textbf{Fine-tuning based approaches} emphasize the effective adaptation of pre-trained LLMs by altering position embeddings to accommodate extended context lengths. Approaches such as Code LLaMA~\cite{codellama}, LLaMA2 Long~\cite{xiong2023effective}, and ScaledRoPE~\cite{liu2023scaling} employ an extremely large base value for RoPE, followed by fine-tuning on their desired sequence length. In contrast, our strategy allows for flexible targeting of multiple context lengths and is capable of handling lengths exceeding 2M tokens. More recently, the high computational requirements of fine-tuning on very long contexts (greater than 128k tokens) have prompted the development of methods like LongLoRA~\cite{longlora} and PoSE~\cite{zhu2023pose}, which seek to lessen this resource burden. Notably, our technique operates independently of these resource-efficient fine-tuning advancements.

\section{Conclusion}

This paper introduces LongRoPE, a technique that achieves an extraordinary extension of LLM context lengths up to 2048k tokens, all the while preserving their proficiency in handling the original, shorter context windows. Our approach leverages two distinct types of non-uniformities inherent in RoPE positional embedding, uncovered through an efficient evolutionary search mechanism. This methodology yields dual advantages: it enables robust initialization for subsequent fine-tuning and permits an 8$\times$ expansion of the context window even without the need for fine-tuning. Furthermore, we outline a staged expansion protocol, wherein LLMs fine-tuned on 256k contexts are incrementally scaled to support a 2048k context window, obviating the necessity for additional fine-tuning. Through comprehensive experiments, we confirm the strong performance of LongRoPE. We anticipate that LongRoPE-2048k will facilitate a range of novel applications requiring extensive context and foster ongoing exploration in this area.

\section*{Broader Impacts} The research outlined in this paper seeks to contribute to the progression of Machine Learning as a discipline. While our work may entail various societal implications, we do not identify any particular issues that require explicit mention at this time.

	\balance

\end{document}